% =============================================================================
%  E89: ARC-AGI-2 via LLM-Written Python — Arbitrary Code, No DSL
%  Technical note: verified-code synthesis without a fixed DSL on ARC-AGI-2.
%  Companion to E84 (T358). Built from the Task Force house template.
% =============================================================================
\documentclass[11pt]{article}
\usepackage[utf8]{inputenc}
\usepackage[margin=1in]{geometry}
\usepackage{booktabs}
\usepackage{amsmath}
\usepackage{amssymb}
\usepackage{hyperref}
\usepackage{xcolor}
\usepackage{listings}
\usepackage{multirow}
\usepackage{array}
\definecolor{kwblue}{rgb}{0.10,0.20,0.65}
\definecolor{cmtgreen}{rgb}{0.20,0.45,0.20}
\definecolor{strred}{rgb}{0.62,0.12,0.12}
\lstdefinestyle{housespecpy}{
  basicstyle=\ttfamily\footnotesize,
  backgroundcolor=\color{black!3},
  frame=single, framesep=5pt, rulecolor=\color{black!45},
  columns=fullflexible, breaklines=true, keepspaces=true,
  showstringspaces=false,
  commentstyle=\color{cmtgreen}\itshape,
  stringstyle=\color{strred},
  keywordstyle=\color{kwblue}\bfseries,
  language=Python,
}
\hypersetup{colorlinks=true, linkcolor=blue, citecolor=blue, urlcolor=blue}

\newcommand{\openworld}{\textsc{OpenWorld}}
\newcommand{\monthyeartoday}{%
  \ifcase\month\or January\or February\or March\or April\or May\or June\or
  July\or August\or September\or October\or November\or December\fi
  \space\number\year}

\begin{document}

\thispagestyle{empty}
\noindent
\begin{minipage}[c]{0.65\textwidth}
  \raggedright\textbf{E89 Technical Note} \quad \texttt{exp-e89-arc-agi2-python-synth}
\end{minipage}\hfill
\begin{minipage}[c]{0.35\textwidth}
  \raggedleft\itshape \monthyeartoday
\end{minipage}

\vspace{1em}

\begin{center}
  {\LARGE\bfseries E89: ARC-AGI-2 via LLM-Written Python}\\[0.4em]
  {\large Arbitrary Code, No DSL, Exact-Match Verify Gate}\\[0.8em]
  {\normalsize D.O.H.\ Research Team \quad Quome}\\[0.3em]
  {\small\texttt{exp-e89-arc-agi2-python-synth} branch, \openworld{} repository}
\end{center}

\vspace{0.5em}
\hrule
\vspace{0.5em}

\begin{abstract}
We test whether removing a fixed DSL constraint improves LLM-based program
synthesis on ARC-AGI-2.  E84 (fixed 17-primitive geometric DSL) scored 0\%
and attributed failure to coverage gaps.  E89 removes that constraint
entirely: the synthesis agent writes arbitrary Python with no vocabulary
restrictions, accepted only when programs \emph{exactly} reproduce every
demo pair.  Burst-1 (30 tasks, 3 candidates/task, \texttt{claude-sonnet-4-6})
scores \textbf{0/30 (0\%)}—identical to E84.  We conclude that the binding
bottleneck is not DSL coverage but LLM rule-inference correctness under
few-shot conditions.
\end{abstract}

\vspace{0.5em}
\hrule

%% ============================================================
\section{Motivation}
\label{sec:motivation}

E84 tested a 17-primitive geometric DSL (translate, recolor, tile, etc.)
against ARC-AGI-2 and scored 0\% across six search budgets (10--10{,}000).
Post-hoc analysis identified the bottleneck as \emph{coverage}: the DSL
lacked connected-component, object-identity, and counting primitives required
by the evaluated tasks.  E84's conclusion was explicit: ``scaling the fixed
DSL does not help; the straitjacket is the representation, not the search.''

E89 tests this claim directly.  If coverage is the bottleneck, replacing the
fixed DSL with unrestricted Python should recover accuracy.  If the LLM
cannot reliably derive correct transformation rules from a handful of
examples regardless of representation, both approaches will score zero.

%% ============================================================
\section{Method}
\label{sec:method}

\subsection{Pipeline}

\begin{enumerate}
  \item \textbf{Load task}: ARC-AGI-2 evaluation task with demo pairs and a held-out test grid.
  \item \textbf{Synthesize}: Send demo pairs to \texttt{claude-sonnet-4-6} via DM protocol
        (arc\_synthesizer\_harness.py).  The prompt asks the model to write a
        \texttt{transform(grid) -> grid} function in \emph{arbitrary Python}
        with no vocabulary restrictions.
  \item \textbf{Verify}: Accept candidates that \emph{exact-match} every demo pair.
        Partial credit is not allowed; any incorrect cell = reject.
  \item \textbf{Apply}: Run verified survivors on the test grid.
        Vote across multiple verified candidates if more than one survives.
\end{enumerate}

\subsection{Configuration}

\begin{center}
\begin{tabular}{ll}
\toprule
Parameter & Value \\
\midrule
Synthesis model & \texttt{claude-sonnet-4-6} \\
Candidates per task & 3 (burst-1) \\
Tasks evaluated & 30 (ARC-AGI-2 evaluation set) \\
Verify gate & Exact match on all demo pairs \\
Evaluation set & Public ARC-AGI-2 evaluation \\
\bottomrule
\end{tabular}
\end{center}

\subsection{Comparison baseline}

E84 (T358): fixed 17-primitive geometric DSL on the same ARC-AGI-2 evaluation set.
Both experiments use identical task selection and verify gate semantics.

%% ============================================================
\section{Results}
\label{sec:results}

\subsection{Burst-1: 30-task evaluation}

\begin{center}
\begin{tabular}{lrr}
\toprule
Metric & E84 (Fixed DSL) & E89 (Arbitrary Python) \\
\midrule
Tasks evaluated & 30 & 30 \\
Candidates per task & 6-budgets & 3 \\
Candidates passing verify gate & 0 & 1 / 90 (1.1\%) \\
Tasks solved (ground truth) & 0 / 30 & \textbf{0 / 30} \\
ARC-AGI-2 accuracy & 0\% & \textbf{0\%} \\
\bottomrule
\end{tabular}
\end{center}

The one candidate that passed the exact-match verify gate did not match the
ground-truth test output.

\subsection{Best-effort verify-gate demonstration (T365)}

To confirm the verify gate is functional end-to-end, a best-effort sub-run
targeted five hand-selected tasks with 32 candidates each.  Task
\texttt{d35bdbdc} produced 20/32 candidates passing all demo pairs via exact
match, with consistent outputs on all three test inputs.  The transformation
pattern (``chain pairing'': shapes with uniform 3$\times$3 surroundings form
chains; alternate shapes swap centers) was decoded correctly.
This demonstrates the pipeline can produce and accept verified programs; the
challenge is generalization across the broad evaluation distribution.

%% ============================================================
\section{Analysis}
\label{sec:analysis}

\subsection{Primary finding: coverage is not the bottleneck}

E84 hypothesized that its 0\% score reflected DSL coverage gaps (no
object-identity, connected-component, or counting primitives).  E89 removes
those restrictions entirely and scores identically.  The data falsify the
coverage-gap hypothesis as the \emph{primary} bottleneck.

\subsection{Alternative hypothesis: inference correctness}

The 1.1\% gate-pass rate (1/90 candidates) across 30 tasks indicates that
\texttt{claude-sonnet-4-6} rarely produces programs that correctly replicate
all demo pairs, regardless of representation.  Candidate programs showed
plausible but incorrect approaches: typically identifying the correct
structural pattern class while mishandling edge cases, directionality, or
selection logic.  The exact-match gate is unforgiving; programs that
``approximately'' capture the pattern still fail.

\subsection{The regime map}

Together, E84 and E89 establish two axes of a synthesis regime map:

\begin{center}
\begin{tabular}{lccc}
\toprule
Experiment & Representation & Bottleneck tested & ARC-AGI-2 \\
\midrule
E84 & Fixed 17-prim DSL & Coverage & 0\% \\
E89 & Arbitrary Python & Inference correctness & 0\% \\
\bottomrule
\end{tabular}
\end{center}

Future axes: multi-turn refinement with failure feedback (E84/E89 both
used single-shot synthesis); stronger models (Opus 4.8); and higher candidate
counts (N=32 expected gate-pass rate $\approx 28$\% if rate is stable at
1.1\%).

%% ============================================================
\section{Limitations}
\label{sec:limitations}

\begin{itemize}
  \item \textbf{Sample size}: 30 tasks from the public evaluation set.
  \item \textbf{Candidates}: 3 per task.  Higher N could recover some tasks.
  \item \textbf{Contamination}: Tasks are from the public set; contamination
        cannot be ruled out, though ARC-AGI-2 is purpose-built to resist
        memorization.
  \item \textbf{No private-set score}: The definitive headline number requires
        the held-out private evaluation set.
\end{itemize}

%% ============================================================
\section{Conclusion}
\label{sec:conclusion}

Removing the DSL constraint does not improve ARC-AGI-2 accuracy.  E89
(arbitrary Python) scores 0/30, matching E84 (fixed geometric DSL).
The binding bottleneck for few-shot program synthesis on ARC-AGI-2 is
LLM rule-inference correctness, not representation coverage.

This is a useful negative result: it rules out one class of explanation
(vocabulary poverty) and points toward a distinct engineering lever
(inference correctness via refinement, stronger models, or larger candidate
pools) as the productive path forward.

%% ============================================================
\appendix
\section{Failure mode examples}
\label{sec:appendix}

Candidate programs in burst-1 consistently showed the following pattern:

\begin{lstlisting}[style=housespecpy]
def transform(grid):
    # Model identifies pattern: "fill enclosed regions"
    # But misses the diagonal-vs-orthogonal distinction
    result = [[row[:] for row in grid]]
    for r in range(1, len(grid)-1):
        for c in range(1, len(grid[0])-1):
            if grid[r][c] == 0 and is_enclosed_ortho(grid, r, c):
                result[r][c] = 1  # Wrong: should check diagonal too
    return result
\end{lstlisting}

The model identifies the correct transformation class (enclosed-region fill)
but implements it incorrectly (orthogonal-only vs.\ true enclosure).  The
exact-match gate correctly rejects these; no partial credit is awarded.

\end{document}
